\chapter{实现补充与复现实验说明}\label{appendix:code}

\section{真实实现说明}

附录不再给出与仓库实现不一致的示例代码，而是摘录当前仓库中的真实核心实现，并用少量注释说明其与正文算法描述的对应关系。本文涉及的核心源码主要位于 \path{src/rbsog_md/forces/rbsog.py}、\path{src/rbsog_md/forces/neighbor.py}、\path{src/rbsog_md/forces/sog.py} 与 \path{src/rbsog_md/benchmark.py}。

当前实现有三个与正文结论直接相关的特征：

（1）短程库仑相互作用由Verlet近邻表求解器负责，支持Numba加速；

（2）长程部分不是简单地“抽取一个固定批量内的所有粒子对”，而是对粒子索引按重要性分布进行采样，再使用重要性权重进行无偏校正；

（3）SOG核通过最小二乘拟合得到，并在给定盒子尺寸下自动重建。

\section{RBSOG求解器真实核心代码节选}

下面节选自 \path{src/rbsog_md/forces/rbsog.py}。这一实现展示了：先计算短程结果，再对长程部分进行重要性采样，最后将两部分贡献合并。

\begin{lstlisting}[language=Python, caption=RBSOG求解器真实核心代码节选]
@dataclass
class RBSOGConfig:
    batch_size: int = 100
    cutoff_short: float = 1.2
    sog_terms: int = 12
    r_min: float = 0.15
    importance_exponent: float = 1.0
    neighbor_skin: float = 0.3
    neighbor_rebuild_interval: int = 10
    use_numba: bool = True


class RBSOGSolver:
    name = "rbsog"

    def __init__(self, config: RBSOGConfig) -> None:
        self.config = config
        self.short_range_solver = NeighborListCoulombSolver(
            cutoff=config.cutoff_short,
            skin=config.neighbor_skin,
            rebuild_interval=config.neighbor_rebuild_interval,
            use_numba=config.use_numba,
        )
        self.long_range_jit_enabled = bool(config.use_numba and _NUMBA_AVAILABLE)
        self.kernel: SOGKernel | None = None
        self.kernel_box: np.ndarray | None = None

    def _ensure_kernel(self, box: np.ndarray) -> None:
        if self.kernel is not None and self.kernel_box is not None and np.allclose(box, self.kernel_box):
            return
        r_max = 0.5 * float(np.min(box))
        self.kernel = fit_sog_kernel(
            r_min=self.config.r_min,
            r_max=r_max,
            n_terms=self.config.sog_terms,
        )
        self.kernel_box = box.copy()

    def compute(self, system: ParticleSystem, rng: np.random.Generator | None = None) -> ForceResult:
        if rng is None:
            rng = np.random.default_rng()

        self._ensure_kernel(system.box)
        assert self.kernel is not None

        short_result = self.short_range_solver.compute(system)
        forces = short_result.forces.copy()
        potential = short_result.potential
        virial = short_result.virial

        positions = system.positions
        charges = system.charges
        n_particles = system.n_particles
        m_samples = max(self.config.batch_size * n_particles, 1)

        q_abs = np.abs(charges) + 1e-6
        proposal = q_abs**self.config.importance_exponent
        proposal /= np.sum(proposal)

        i_all = rng.choice(n_particles, size=m_samples, p=proposal)
        j_all = rng.choice(n_particles, size=m_samples, p=proposal)
        mask_not_self = i_all != j_all

        i_idx = i_all[mask_not_self]
        j_idx = j_all[mask_not_self]
        displacement = minimum_image(positions[j_idx] - positions[i_idx], system.box)
        r_sq = np.einsum("ij,ij->i", displacement, displacement)

        cutoff_sq = self.config.cutoff_short * self.config.cutoff_short
        r_min_sq = self.config.r_min * self.config.r_min
        mask_long = (r_sq > cutoff_sq) & (r_sq > r_min_sq)

        i_valid = i_idx[mask_long]
        j_valid = j_idx[mask_long]
        disp_valid = displacement[mask_long]
        r_sq_valid = r_sq[mask_long]

        pair_prob = proposal[i_valid] * proposal[j_valid]
        sample_weights = 1.0 / (m_samples * pair_prob)

        pref = self.kernel.force_prefactor(r_sq_valid)
        qij = charges[i_valid] * charges[j_valid]
        pair_force = (qij * pref)[:, None] * disp_valid

        scaled_force = sample_weights[:, None] * pair_force
        np.add.at(forces, i_valid, scaled_force)
        np.add.at(forces, j_valid, -scaled_force)

        pair_potential = qij * self.kernel.evaluate_from_r_sq(r_sq_valid)
        potential += float(np.sum(0.5 * sample_weights * pair_potential))
        virial += float(np.sum(0.5 * sample_weights * np.einsum("ij,ij->i", disp_valid, pair_force)))

        return ForceResult(forces=forces, potential=potential, virial=virial)
\end{lstlisting}

这一节选说明，正文中“随机批量采样”的实现应理解为“基于粒子电荷幅值的成对重要性采样”，而不是简单的均匀批量抽样。

\section{近邻表与短程求解真实实现节选}

下面代码节选自 \path{src/rbsog_md/forces/neighbor.py}，展示了Verlet近邻表的重建判据和短程库仑求解流程。

\begin{lstlisting}[language=Python, caption=近邻表与短程求解真实代码节选]
@dataclass
class VerletNeighborList:
    cutoff: float
    skin: float = 0.3
    rebuild_interval: int = 10

    @property
    def pair_cutoff(self) -> float:
        return self.cutoff + self.skin

    def _needs_rebuild(self, system: ParticleSystem) -> bool:
        if self._pairs is None or self._ref_positions is None or self._ref_box is None:
            return True
        if not np.allclose(system.box, self._ref_box):
            return True
        if self._steps_since_build >= self.rebuild_interval:
            return True
        if self.skin <= 0.0:
            return False

        displacement = minimum_image(system.positions - self._ref_positions, system.box)
        max_disp_sq = float(np.max(np.einsum("ij,ij->i", displacement, displacement)))
        return max_disp_sq > (0.5 * self.skin) ** 2


class NeighborListCoulombSolver:
    def compute(self, system: ParticleSystem, rng: np.random.Generator | None = None) -> ForceResult:
        del rng
        pairs = self.neighbor_list.get_pairs(system)
        if pairs.size == 0:
            return ForceResult(forces=np.zeros_like(system.positions), potential=0.0, virial=0.0)

        return self._compute_vectorized(
            positions=system.positions,
            charges=system.charges,
            box=system.box,
            pairs=pairs,
            cutoff_sq=self.cutoff * self.cutoff,
            softening_sq=self.softening * self.softening,
        )
\end{lstlisting}

\section{SOG核拟合真实实现节选}

下面代码节选自仓库中的SOG核拟合实现文件：

\noindent\path{src/rbsog_md/forces/sog.py}

其核心是：在给定$[r_{\min}, r_{\max}]$区间内，对$1/r$进行最小二乘拟合，得到用于长程近似的高斯核参数。

\begin{lstlisting}[language=Python, caption=SOG核拟合真实代码节选]
def fit_sog_kernel(
    r_min: float,
    r_max: float,
    n_terms: int = 12,
    n_samples: int = 4096,
) -> SOGKernel:
    alpha_min = 0.5 / (r_max * r_max)
    alpha_max = 8.0 / (r_min * r_min)
    alphas = np.geomspace(alpha_min, alpha_max, n_terms)

    r = np.geomspace(r_min, r_max, n_samples)
    design = np.exp(-np.outer(r * r, alphas))
    target = 1.0 / r

    weights, _, _, _ = np.linalg.lstsq(design, target, rcond=None)
    weights = np.clip(weights, 0.0, None)

    return SOGKernel(alphas=alphas, weights=weights, r_min=r_min, r_max=r_max)
\end{lstlisting}

\section{复现实验命令}

本文最终采用的实验流程与仓库中的命令行入口一致。对应协议可见 \path{docs/experiment_protocol.md}，核心命令如下：

\begin{lstlisting}[language=bash, caption=论文核心实验命令]
# PPPM vs RBSOG 主基准
rbsog-md benchmark \
  --solvers pppm,rbsog \
  --steps 1000 \
  --sample-interval 20 \
  --seeds 3 \
  --n-lipids 128 \
  --n-solvent 512 \
  --batch-size 100 \
  --grid 32 32 32 \
  --out results/thesis_final_benchmark_main

# 批量大小扫描
rbsog-md sweep-batch \
  --steps 1000 \
  --sample-interval 20 \
  --seeds 3 \
  --batch-sizes 50,100,200,400 \
  --objective-time-weight 2 \
  --objective-variance-weight 1 \
  --out results/thesis_final_batch_sweep

# JIT 消融
rbsog-md benchmark \
  --solvers rbsog \
  --steps 400 \
  --sample-interval 10 \
  --seeds 2 \
  --batch-size 100 \
  --out results/thesis_ablation_jit_on

rbsog-md benchmark \
  --solvers rbsog \
  --steps 400 \
  --sample-interval 10 \
  --seeds 2 \
  --batch-size 100 \
  --disable-numba \
  --out results/thesis_ablation_jit_off
\end{lstlisting}

\section{附录小结}

附录的目的不是重复展示与正文无关的大段模板代码，而是说明论文中对RBSOG、近邻表、SOG核拟合以及实验流程的描述都能在当前仓库真实实现中找到对应依据。这样可以确保正文、附录、源码三者之间保持一致。
